\documentclass{article}

% if you need to pass options to natbib, use, e.g.:
\PassOptionsToPackage{numbers, compress}{natbib}
% before loading neurips_2019

% ready for submission
% \usepackage{neurips_2019}

% to compile a preprint version, e.g., for submission to arXiv, add add the
% [preprint] option:
% \usepackage[preprint]{neurips_2019}

% to compile a camera-ready version, add the [final] option, e.g.:
\usepackage[final]{neurips_2019}

% to avoid loading the natbib package, add option nonatbib:
%     \usepackage[nonatbib]{neurips_2019}

\usepackage[utf8]{inputenc} % allow utf-8 input
\usepackage[T1]{fontenc}    % use 8-bit T1 fonts
\usepackage{hyperref}       % hyperlinks
\usepackage{url}            % simple URL typesetting
\usepackage{booktabs}       % professional-quality tables
\usepackage{amsfonts}       % blackboard math symbols
\usepackage{nicefrac}       % compact symbols for 1/2, etc.
\usepackage{microtype}      % microtypography

\usepackage{mathtools}
\usepackage{tabularx}
\usepackage{amsmath}
\usepackage{amssymb}
\usepackage[noend]{algpseudocode}
\usepackage{algorithm,algorithmicx}
\usepackage{graphicx}
\usepackage{color}
\usepackage{subfig}
\usepackage{caption}
\usepackage{array}
\usepackage{wrapfig}
\usepackage[final]{pdfpages}

\captionsetup[table]{skip=5pt}

\newcommand*\Let[2]{\State #1 $\gets$ #2}
\algrenewcommand\algorithmicrequire{\textbf{Precondition:}}
\algrenewcommand\algorithmicensure{\textbf{Postcondition:}}

\newcommand{\ra}[1]{\renewcommand{\arraystretch}{#1}}
\newcolumntype{F}{>{\centering\arraybackslash}p{1em}}
\newcolumntype{C}{>{\centering\arraybackslash}p{6em}}


\graphicspath{{./figures/}}

%\DeclareMathOperator*{\argmin}{arg\,min}
\DeclareMathOperator*{\argmin}{arg\,min}
\DeclareMathOperator*{\intr}{\mathbf{K}}
\DeclareMathOperator*{\extr}{\mathbf{E}}

\newcommand{\OURS}{Scene Representation Networks}
\newcommand{\OURSHORT}{SRNs}

\title{\OURS: Continuous 3D-Structure-Aware Neural Scene Representations}

\newcommand{\VS}[1]{#1}
%\newcommand{\VS}[1]{\textcolor{red}{#1}}
\newcommand{\GW}[1]{\textcolor{blue}{#1}}
\newcommand{\MZ}[1]{\textcolor{green}{#1}}

% The \author macro works with any number of authors. There are two commands
% used to separate the names and addresses of multiple authors: \And and \AND.
%
% Using \And between authors leaves it to LaTeX to determine where to break the
% lines. Using \AND forces a line break at that point. So, if LaTeX puts 3 of 4
% authors names on the first line, and the last on the second line, try using
% \AND instead of \And before the third author name.

\author{
	\begin{tabular}{c@{\hspace*{.6cm}}c@{\hspace*{.6cm}}c@{\hspace*{.35cm}}}
		Vincent Sitzmann & Michael Zollhöfer & Gordon Wetzstein\\
	\end{tabular}\\[0.1cm]
	\texttt{\{sitzmann,\:zollhoefer\}@cs.stanford.edu, gordon.wetzstein@stanford.edu}\\
	Stanford University\\
}



\begin{document}

\maketitle

	\vbox{%
	\vskip -0.26in
	\hskip -0.15in
	\hsize\textwidth
	\linewidth\hsize
	\centering
	\tt\normalsize\href{https://vsitzmann.github.io/srns/}{vsitzmann.github.io/srns/}
	\vskip 0.3in
}

\begin{abstract}
\VS{
Unsupervised learning with generative models has the potential of discovering rich representations of 3D scenes.
%
While geometric deep learning has explored 3D-structure-aware representations of scene geometry, these models typically require explicit 3D supervision.
%
Emerging neural scene representations can be trained only with posed 2D images, but existing methods ignore the three-dimensional structure of scenes.}
%
We propose Scene Representation Networks (SRNs), a continuous, 3D-structure-aware scene representation that encodes both geometry and appearance. SRNs represent scenes as continuous functions that map world coordinates to a feature representation of local scene properties. By formulating the image formation as a differentiable ray-marching algorithm, SRNs can be trained end-to-end from only 2D images and their camera poses, without access to depth or shape. This formulation naturally generalizes across scenes, learning powerful geometry and appearance priors in the process. We demonstrate the potential of SRNs by evaluating them for novel view synthesis, few-shot reconstruction, joint shape and appearance interpolation, and unsupervised discovery of a non-rigid face model.\footnote{Please see supplemental video for additional results.}
\end{abstract}

\section{Introduction}
\label{sec:intro}
\input{sections/introduction}

\section{Related Work}
\label{sec:related}
\input{sections/related}

\section{Formulation}
\label{sec:formulation}
\input{sections/formulation}

\section{Experiments}
\label{sec:experiments}
\input{sections/experiments}

\section{Discussion}
\label{sec:discussion}
\input{sections/discussion}

\section{Acknowledgements}
\label{sec:acks}
\input{sections/acks}

{\small
	\bibliographystyle{IEEEtranN}
	\bibliography{references}
}

\includepdf[pages=-]{SRNs_supp.pdf}

\end{document}
